%!TEX encoding = UTF8
%!TEX root = 0-notes.tex

\chapter*{Préface}


\section*{Syllabus}

\subsection*{Description}

\begin{itemize}
	\item
	Panorama des structures algébriques : groupes, anneaux, corps, algèbres avec des exemples de L1.
	\item
	L'algèbre $\K[x]$ : définition, opérations degré, $\K_n[x]$ ($\K = \Q, \R$ ou $\C$).
	\item
	Arithmétique de $\K[x]$ :
	\begin{itemize}
		\item
		divisibilité, polynômes irréductibles, division euclidienne, algorithme d'Euclide, PGCD et PPCM, théorème de Bézout,
		lemme de Gauss, décomposition en facteurs irréductibles.
		\item
		Notion d'idéal d'un anneau, $\Z$ et $\K[x]$ comme anneaux principaux, ré-interprétation de divisibilité, PGCD, PPCM, en termes d'idéaux.
	\end{itemize}
	\item
	Fonctions polynomiales :
	\begin{itemize}
		\item
		rappels : racines, multiplicité, dérivation, formule de Taylor, caractérisation de la multiplicité des racines.
		\item
		polynôme scindé, relation racines-coefficients, Théorème de D'Alembert-Gauss, décomposition en facteurs irréductibles dans $\R[x]$ et $\C[x]$. 
		Racines $n$-ièmes de l'unité.
	\end{itemize}
	\item
	Fractions rationnels : définition comme corps des fractions de $\K[x]$.
	Degré, partie entière, décomposition en éléments simples (sur $\R$ et $\C$).
\end{itemize}

\subsection*{Objectifs}

Dans ce cours, on introduira un panorama des structures algébriques (anneau, idéal, corps) avant d'aborder l'algèbre $\K[x]$ et définir l'arithmétique
sur les polynômes en faisant des parallèles avec l'arithmétique des entiers vue en L1.
Des partie calculatoires sur les fonctions polynomiales et les fractions rationnelles seront traitées (factorisation, décompositions explicites).

\subsection*{Pré-requis nécessaires}

MG203 - Arithmétique et dénombrement

\subsection*{Pré-requis recommandés}

L1 maths.

\section*{Structure du cours}

Construction d'un cours de mathématiques : Exemple -- Remarque -- Définition -- Lemme -- Proposition -- Théorème -- Démonstration -- Corollaire.

Voir \emph{Mechanica, volume 1}, Leonhard Euler, 1736 (\href{https://scholarlycommons.pacific.edu/euler-works/15/}{E15}), pour un exemple de cours de mathématiques en latin.

Wolframalpha permet de faire du calcul exact.

\begin{definition*}
Associe un mot de vocabulaire à un objet mathématique.
Un mot de vocabulaire peut avoir un sens différent en mathématiques qu'en langage vernaculaire.
C'est par exemple le cas de « ensemble », « couple », « rationnel », « réel », ...
\end{definition*}

\begin{lemme*}
	Résultat préliminaire, parfois technique, utile à la démonstration d'une proposition plus importante.
\end{lemme*}

\begin{proposition*}
	Énoncé vrai et important.
\end{proposition*}

\begin{theorem*}
	Énoncé vrai et fondamental.
\end{theorem*}

\begin{corollaire*}
	Résultat particulier, découlant souvent quasi directement d'un théorème ou d'une proposition.
\end{corollaire*}

\newcounter{preface}
\begin{Exercise}[label=1, counter=preface]
	Exemple d'exercice d'application directe.
\end{Exercise}
\begin{Exercise}[difficulty=1, label=2, counter=preface]
	Exemple de problème.
\end{Exercise}
\begin{Exercise}[difficulty=2, label=3, counter=preface]
	Exemple de problème d'approfondissement.
\end{Exercise}
\shipoutExercise
\begin{Answer}[ref=1]
	Solution 1.
\end{Answer}
\begin{Answer}[ref=2]
	Solution 2.
\end{Answer}
\begin{Answer}[ref=3]
	Solution 3.
\end{Answer}
\begin{multicols}{3}
\shipoutAnswer
\end{multicols}

%\setcounter{Exercise}{0}
